\section{A concrete path to complete Lean formalization}
\label{sec:formalization}

\subsection{Three distinct levels of completion}

The mathematical proofs in this paper are complete human-readable
arguments, with their computation and induction bounds given explicitly.
The supplied Lean files are an implementation blueprint. They have not
been compiled in this environment, and this package contains no finished
Lean term proving the four base cases or their internal $\ipra$
derivations. This distinction is part of the deliverable, not an implicit
assumption about a future formalization.

There are three levels to keep separate:
\begin{enumerate}
\item \textbf{Arithmetic semantics.} Prove the displayed integer identities,
number-theoretic lemmas, and FLT assertions in Lean's ordinary logic.
\item \textbf{Programs and restricted induction.} Represent every numerical
operation by an explicit primitive recursive code, and identify every
induction formula as quantifier free. A termination proof alone is not
enough to establish this level.
\item \textbf{An internal derivation.} Construct an inhabitant of
\texttt{Prf PRAAx [] sentence} in ProveIt's actual syntax, using only
the theory's defining equations and allowed rules.
\end{enumerate}
An ordinary Lean theorem at level 1 does not automatically give level 3.
Even an ambient proof using no call to \texttt{classical} would not establish
that its induction is restricted to the induction scheme of $\ipra$.

\subsection{The exact repository target}

The inspected ProveIt snapshot is
\begin{center}\small\ttfamily
bb4427b10219ad61e0d2875fdde04bcd23b47ec6.
\end{center}
Its root toolchain is Lean \texttt{v4.32.0}, and its Mathlib dependency is
\texttt{v4.32.0}. The relevant modules are:
\begin{center}
\begin{tabular}{>{\raggedright\arraybackslash}p{0.31\textwidth}>{\raggedright\arraybackslash}p{0.61\textwidth}}
\toprule
Module under \texttt{ConciseIPRA} & Role verified in the source\\
\midrule
\texttt{Syntax} & Scoped terms, formulas, substitution, and the predicate
\texttt{Fm.QF}.\\
\texttt{Deduction} & Intuitionistic derivations \texttt{Prf}, including
quantifier rules and equality substitution.\\
\texttt{PR} & Function codes \texttt{PR}, terms \texttt{PTm}, the theory
\texttt{PRAAx}, and the rule \texttt{PRA.qfInd}.\\
\texttt{NativeBiInterpretation} & The equivalence
\texttt{pra\_iff\_native}, which translates a supplied internal proof.\\
\texttt{TeqRec} & A derived primitive recursor in the iterator language,
useful if implementing directly in the scalar core.\\
\bottomrule
\end{tabular}
\end{center}
The snapshot and the inspected blob hashes are recorded in
\texttt{evidence/repository\_snapshot.json}. No result about FLT has been
assumed merely because these foundational modules are present.

The file \texttt{lean/ObjectLanguageTargets.lean} defines explicit codes
for addition, multiplication, and exponentiation using
\texttt{PR.comp} and \texttt{PR.prec}. The convention is
\[
 \operatorname{addCode}(n,x)=n+x,\quad
 \operatorname{mulCode}(n,x)=nx,\quad
 \operatorname{powCode}(n,x)=x^n;
\]
the recursion argument is always first, as required by the repository.
It then builds the actual formula
\[
 \operatorname{fixedSentence}(d)=
 \forall x\,\forall y\,\forall z\,
 \bigl(x\ne0\to y\ne0\to z\ne0\to
            x^d+y^d=z^d\to\bot\bigr)
\]
using \texttt{Fm.eq}, \texttt{Fm.imp}, \texttt{Fm.bot}, and
\texttt{Fm.all}. Its definition
\[
 \texttt{InternalFLT d}
 :=\texttt{Prf PRAAx [] (fixedSentence d)}
\]
is the precise type of the missing formal proof object. It is not a
synonym for an ambient Lean proposition about natural numbers.

The uniform formula uses a free exponent, a nonzero guard, and a
disjunction of the four formulas $\exists k\,dk=n$. A nonzero multiple of
one of $3,4,5,7$ is automatically at least $3$, so this is equivalent to
the stated main family. The bound for selecting $k$ is $k\leq n$.
The code also gives the exact conditional translation to
\texttt{NPrf [] (NFm.ofPRA sentence)} using
\texttt{pra\_iff\_native}. This translation assumes an actual internal
derivation as its input; it supplies none of the arithmetic work.

\subsection{Required numerical definitions}

The file \texttt{lean/ArithmeticBlueprint.lean} records the semantic
definitions. The internal version should implement the same operations
on numerical codes, in the following order.

\paragraph{Natural arithmetic.}
Define truncated subtraction, the zero test, comparison, quotient and
remainder with a prescribed value at divisor zero, bounded sums and
products, and a first-success bounded selector. For a predicate
$\chi(i,\vec a)$, the selector returns the least $i\leq B$ with
$\chi(i,\vec a)=1$, together with a success flag. Its invariant, after
scanning $0,\ldots,j$, states that the returned candidate is either the
least success in that interval or the flag is zero and all those tests
are zero. The latter bounded universal is encoded by a product. The
invariant is therefore quantifier free after the definitions are expanded.

\paragraph{Signed arithmetic and tuples.}
Use the signed code and pairing function from
Section~\ref{sec:foundations}; implement addition, multiplication, negation,
absolute value, and sign by finite case splits. Prove that the explicit
operations satisfy the ring identities. A code for an integer is not a
quantifier over a new infinite sort: every statement expands to a numerical
formula. Rational quadratic coordinates used in the quintic proof can be
stored as integer numerators with positive integer denominators, and all
comparisons can be made after cross multiplication.

\paragraph{Divisibility and roots.}
Define the gcd computation with its clock, its Bezout accumulator, prime
testing, valuation at nonzero input, finite factor lists, and an exact-root
selector. For exponent $e\geq1$ the root selector only scans $0,\ldots,N$.
Its correctness condition is a finite equality test. Its success in the
proof follows from the coprime-factorization lemma, not from a numerical
experiment. On the few invalid inputs where a valuation is not defined
mathematically, the code returns a specified default, and every theorem
has an explicit nonzero guard.

\paragraph{The two quadratic pair rings.}
Define multiplication, conjugation, norm, rational quotient coordinates,
nearest-integer rounding, and a clocked gcd for each pair ring. Prove norm
multiplicativity by expanding polynomial expressions. The Eisenstein norm
is positive definite, so a norm bound directly bounds coordinates. The
golden norm is indefinite, so one must first implement the bounded unit
normalization from Section~\ref{sec:five}. Omitting this difference would
invalidate the bounded-factorization claim for the quintic case.

\subsection{State predicates and the central correctness contracts}

Here are the complete mathematical contents of the main state predicates.
They should be Boolean functions on encoded tuples; the corresponding
Lean predicates in the blueprint are readable specifications.

\begin{center}
\small
\begin{tabular}{>{\raggedright\arraybackslash}p{0.14\textwidth}>{\raggedright\arraybackslash}p{0.65\textwidth}>{\raggedright\arraybackslash}p{0.11\textwidth}}
\toprule
State & Predicate, in addition to the displayed equation & Rank\\
\midrule
Cubic $(x,y,z)$ & Nonzero signed pairwise coprime entries;
$x^3+y^3+z^3=0$; $3\mid z$, $3\nmid xy$. & $|z|$\\
Quartic $(x,y,z)$ & Positive entries; $\gcd(x,y)=1$;
$x^4+y^4=z^2$. & $z$\\
Quintic odd $(h,t,s,w)$ & $h\geq0$; $t,s,w>0$;
$\gcd(t,s)=1$; $t,s$ odd; $5\mid s$, $5\nmid t$;
$5^h sF(t,s)=16w^5$. & $s$\\
Quintic even $(h,t,s,w)$ & Same positivity and coprimality;
$t$ odd, $s$ even; $5\mid s$, $5\nmid t$;
$2\cdot5^h sF(t,s)=w^5$. & $s$\\
Septimal $(a,P,Q,R)$ & $a\geq1$; positive odd pairwise coprime
$P,Q,R$; equation \eqref{eq:seven-quartic}. & $a$\\
\bottomrule
\end{tabular}
\end{center}

For each of the first four rows, construct a total selector $D$ and prove
\[
 \chi(s)=1\quad\Longrightarrow\quad
       \chi(D(s))=1\ \wedge\ r(D(s))+1\leq r(s).
\tag{C}
\]
The object-level proof of (C) consists exactly of the corresponding
descent section. Invalid states return a default and do not enter its
hypotheses. The final contradiction is the shared quantifier-free
invariant (I) from Lemma~\ref{lem:clocked}.

For the septimal row one may also use the simpler fixed-box Boolean
\[
 \operatorname{bad7}(a,M)=
 \max_{1\leq P,Q,R\leq M}\chi_{D_a}(P,Q,R).
\]
Here ``max'' means an explicitly nested finite Boolean fold. Prove
$\operatorname{bad7}(1,M)=0$ and
$\operatorname{bad7}(a+1,M)=0$ from
$\operatorname{bad7}(a,M)=0$ for $a\geq1$. The fact that each new
coordinate is at most the old $R$ keeps the same $M$ throughout.
The induction formula is the literal equation
$\operatorname{bad7}(a+1,M)=0$.

In the natural-number implementation of the quartic polynomial, avoid
truncated subtraction. Use the equivalent equality
\[
 P^2+3\cdot7^4\,2^{2a}Q^2R^2
     =Q^4+7^7\,2^{4a+4}R^4.
\]
Intermediate square completions can use the signed-number code. This
small representation decision prevents the formal statement from
silently becoming a different equation.

\subsection{Construction of internal proofs, in dependency order}

The following is a finite implementation plan with identified proof
obligations. It does not call for any new number-theoretic theorem beyond
the arguments in this paper.

\begin{enumerate}
\item \textbf{Arithmetic equations.} Build internal proofs for addition,
multiplication, powers, order, and signed operations. Prove the identities
used by polynomial normalization. At this stage a normalizer may compute
the two expanded polynomials, but it must also emit a derivation using
these laws and \texttt{Prf.eqSubst}. Merely running the ambient Lean
\texttt{ring} tactic does not emit an $\ipra$ proof automatically.

\item \textbf{Bounded-search reflection.} Prove internally that the Boolean
folds and first-success selector satisfy their specified equations and
produce witnesses within their bounds. This is repeated use of
\texttt{PRA.precZero}, \texttt{PRA.precSucc}, and \texttt{PRA.qfInd}.
The invariants described above are explicit recipes for the required
\texttt{Fm.QF} proofs.

\item \textbf{Integer arithmetic tools.} Formalize
Lemma~\ref{lem:integer-tools} and primitive reduction. Use bounded lists
of prime factors; replace any induction over ``all smaller integers''
by a Boolean invariant over the finite list or bounded range. Prove the
valuation rule for a sum of unequal valuations by extracting the smaller
power and reducing the remaining equality modulo the prime.

\item \textbf{Quadratic arithmetic.} Formalize the coordinate identities,
Euclidean gcd, primality of irreducibles, finite factorization, and units.
For $E$ prove the two exact $\lambda$ divisibility tests and the six-unit
table. For $R$ prove bounded unit normalization, the classification
$\pm\varphi^j$, and the modulo-$5$ fifth-power scalar test. These are
the precise algebraic inputs to the root-extraction routines.

\item \textbf{Local transitions.} Prove (C) for cubic, quartic, and the two
quintic states. For exponent $7$, prove the base residue contradiction,
the complete four-way factor allocation, the row-III contradiction, and
the transition to parameter $a-1$. All finite allocations are to be
proved exhaustive before any branch is discharged.

\item \textbf{Finite termination.} Apply the single generic clocked-descent
derivation, or the septimal fixed-box induction. This yields the four
impossibility lemmas without an unbounded minimality predicate.

\item \textbf{Initial transformations.} Translate an arbitrary alleged
Fermat solution to a valid state. Use the explicit gcd reduction,
modulo-$9$ or modulo-$25$ calculations, two quintic parity branches,
and the septimal symmetric-polynomial and valuation argument. All
divisions must carry their already proved divisibility witnesses.

\item \textbf{Closure and export.} Construct \texttt{InternalFLT 3},
\texttt{InternalFLT 4}, \texttt{InternalFLT 5}, and \texttt{InternalFLT 7}.
Formalize exponent lifting and the four-way divisibility split to obtain
\texttt{InternalCoverage}. Finally apply \texttt{pra\_iff\_native} if a
proof in ConciseIPRA's native scalar language is desired.
\end{enumerate}

The source file declares a structure listing the required proof objects,
but does not supply an inhabitant. There are no axioms or admitted holes
pretending to fill them. The ambient blueprint similarly makes any use
of unimplemented base proofs an explicit hypothesis. This allows a future
implementation to replace each interface with a checked definition while
preserving an honest inventory of what remains.

\subsection{Build and acceptance criteria}

The \texttt{lean/README.txt} gives the commands for checking the blueprint
against the pinned ProveIt environment. A completed implementation should
be accepted only after all of the following have been checked:
\begin{enumerate}
\item The code compiles at the pinned toolchain, and the four base results
are proved without placeholders or extra axioms.
\item The completed deliverable contains the actual proof objects
\texttt{InternalFLT d}, not only the semantic statements.
\item Every inductive invariant used in those internal proofs has a
checked \texttt{Fm.QF} derivation. All algorithms used as internal
function symbols have actual \texttt{PR} codes.
\item The proof audit lists dependencies of the final declarations.
The internal proof rules must still be exactly \texttt{PRAAx} and
\texttt{Prf}; a powerful ambient fact must not enter by an added
object-language axiom.
\end{enumerate}
These criteria describe how to finish the requested formalization. They
are not represented as checks already completed by this package.

\section{Artifact verification and reproducibility}
\label{sec:verification}

The Python files serve three limited and useful roles: checking algebraic
transcription, checking finite residue tables, and enumerating the exact
exponent family. They are not proof checkers for $\ipra$ and do not replace
the human proof.

\begin{center}
\begin{tabular}{>{\raggedright\arraybackslash}p{0.36\textwidth}>{\raggedright\arraybackslash}p{0.56\textwidth}}
\toprule
Artifact & Verified role\\
\midrule
\texttt{verify\_cubic.py} & Pair-ring identities, 28,561 multiplication
and Euclidean-division samples, and the six-unit coefficient table.\\
\texttt{verify\_quintic.py} & Exact symbolic norm and fifth-power
identities; finite modulo-$2$, modulo-$5$, and modulo-$25$ tables;
bounded Euclidean and unit-normalization samples.\\
\texttt{verify\_seven.py} & Seven exact symbolic identities, including
both correction square completions and the Fermat-to-quartic substitution.\\
\texttt{coverage.py} & Exact $276/420$ residue count, exponent-reduction
witnesses, and a CSV through exponent $200$.\\
\texttt{ObjectLanguageTargets.lean} & Actual function codes and target
formulas for the internal formalization; uncompiled blueprint.\\
\texttt{ArithmeticBlueprint.lean} & Readable arithmetic definitions and
formalization interfaces; uncompiled blueprint.\\
\bottomrule
\end{tabular}
\end{center}

The algebra scripts use exact integers or rational arithmetic; the
symbolic checks use SymPy rather than floating point. Their output is
recorded in \texttt{evidence/verification.txt}. Mathematical review also
checked the unbounded arguments that sampling cannot establish: finite
factorization, the golden-ring unit normalization, all four septimal
factor allocations, the valuation extraction, the strict descent
inequalities, and the quantifier-free induction formulas.

The archive includes this article's complete modular LaTeX source,
the compiled PDF, the Lean blueprint, the verification scripts, the
source-snapshot manifest, and reproducibility instructions. The PDF was
compiled and rendered for layout inspection. The Lean files were not
compiled; there was no Lean executable in the delivery environment.
\texttt{SHA256SUMS.txt} identifies the exact files in the package.

